\chapter{Introduction}
\label{chap:intro}

\scaffold{Open with the landscape: modern CPUs rely on largely undocumented
performance mechanisms---speculative and out-of-order execution, branch
prediction, prefetching, store-to-load forwarding---that have repeatedly turned
into security vulnerabilities. Spectre and Meltdown showed that speculation plus a
microarchitectural side channel leaks across security boundaries; despite years of
mitigations, new transient-execution leaks keep appearing. Draw from the uasc2027
introduction.}

\section{The Arm blind spot}
\label{sec:intro-arm}
\scaffold{Arm powers almost all smartphones and is spreading into laptops, HPC and
the datacenter, yet has received far less microarchitectural-security scrutiny than
x86-64. It is nonetheless exposed to Spectre-class attacks and to Arm-specific ones
that target its own defenses---PACMAN against Pointer Authentication and TikTag
against the Memory Tagging Extension. Motivate why automated, black-box discovery
on Arm is needed.}

\section{Preconditions as a source of leakage}
\label{sec:intro-precond}
\scaffold{State the recurring pattern this thesis builds on: instructions impose
preconditions on their operands---a matching MTE tag, a valid pointer signature, a
canonical/aligned address---and violating a precondition raises a fault or leaves
another attacker-observable microarchitectural effect. Architecturally these never
retire on the well-formed path; speculatively they decide whether a line or a
predictor entry is allocated.}

\section{Contributions}
\label{sec:intro-contributions}
\scaffold{Enumerate, in the thesis's relevance order: (1) a port of \revizor{} (a
black-box model-based relational-testing fuzzer) to AArch64, with headless fuzzing
of remote CPUs and support for Arm-specific features (MTE, PAuth); (2) generalized
priming---a provably stronger leak-detection rule that detects every leak current
priming detects and strictly more, with two optimal localizers; (3) precondition
fuzzing---a methodology for systematically fuzzing instruction preconditions; (4)
the discovery of a previously-unknown canonicality- and alignment-based leak on
loads, stores and indirect branches, demonstrated on the Neoverse~N3 and Google
Tensor~G3, crossing privilege (EL0--EL2) and virtualization (guest-to-host)
boundaries, and weaponized end-to-end (OpenSSL RSA-2048 key recovery).}

\section{Thesis organization}
\label{sec:intro-organization}
\scaffold{One short paragraph mapping the chapters: background
(\autoref{chap:background}), the AArch64 port (\autoref{chap:port}), generalized
priming (\autoref{chap:priming}), precondition fuzzing (\autoref{chap:precond}),
the canonicality leak (\autoref{chap:canon}), evaluation (\autoref{chap:eval}),
related work (\autoref{chap:related}) and conclusions (\autoref{chap:conclusion}).}
